\section{Results \& technical interpretation}

\subsection{Payment Format Risk Analysis}

The PySpark aggregation identified differences in laundering activity across payment
formats. Automated Clearing House (ACH) transactions represented the largest share of laundering transactions, accounting for approximately 73.6\% of all laundering transactions in the processing
dataset. ACH transactions are a way to move money directly from one bank account to another without using paper cheques, cash, or credit cards. The laundering rate for ACH transactions was approximately 0.331\%, compared with an overall dataset laundering rate of approximately 0.055\%, resulting in
a relative risk of approximately 6.01 times the dataset baseline. In contrast, wire
transactions had a laundering rate close to zero, while reinvestment transactions
contained no labelled laundering transactions in the analysed data. The relative laundering risk by payment format can be seen in \autoref{fig:laundering_format}:

\begin{figure}[H]
    \centering
    \includegraphics[width=0.9\textwidth]{figures/laundering_by_format.png}
    \caption{Relative Laundering Risk by Payment Format}
    \label{fig:laundering_format}
\end{figure}

These results show that laundering activity is not distributed uniformly across payment
formats. The groupBy() and aggregation operations in PySpark allow the transaction
population to be grouped by payment format and reduced to transaction counts,
laundering counts, laundering rates and relative risk measures. This provides a scalable
method for comparing transaction characteristics across a large transaction population.

\subsection{Account-pair Frequency}

The account-pair analysis identified a strong relationship between transaction frequency
and laundering rate. The laundering rate decreases as the number of transactions
between an account pair increases. Relative to the overall dataset laundering rate of
approximately 0.055\%, the 0.73\% rate for single-transaction pairs is approximately 13.2
times higher, while pairs with 2–5 transactions have a rate approximately 9.4 times the
baseline. The laundering rate by account-pair frequency can be seen in \autoref{fig:laundering_pairs}:

\begin{figure}[H]
    \centering
    \includegraphics[width=0.9\textwidth]{figures/laundering_by_pairs.png}
    \caption{Laundering Rate by Account-Pair Transaction Frequency}
    \label{fig:laundering_pairs}
\end{figure}

These results show that, within the analysed dataset, laundering activity is more
concentrated among low-frequency account relationships than among highly recurring
relationships. This demonstrates the value of performing relationship-level aggregation rather than
analysing transactions independently. The PySpark process groups transactions
according to the originating and receiving account combination and then calculates the
number of transactions and laundering transactions associated with each account pair.
These account-level results can subsequently be grouped into transaction-frequency
categories to compare laundering rates across different relationship types.

\subsection{Temporal Laundering Patterns}

The daily aggregation examines how laundering activity varies over time. The results
show variation in the observed laundering rate across the analysis period, indicating
that laundering activity is not uniformly distributed across days. This distribution can be seen in \autoref{fig:laundering_daily}:

\begin{figure}[H]
    \centering
    \includegraphics[width=\textwidth]{figures/daily_laundering.png}
    \caption{Daily Money-Laundering Rate}
    \label{fig:laundering_daily}
\end{figure}

PySpark groups transactions by the derived \emph{Date} field and calculates the total
transaction count and laundering count for each day, from which the daily laundering
rate is calculated. The resulting time-series shows that laundering activity varies across
the analysis period rather than being uniformly distributed. The analysis also indicates
higher observed laundering activity during periods of lower overall transaction volume,
consistent with the Part 1 finding that laundering occurred mainly during weekends. This
demonstrates how distributed grouping and aggregation can transform transaction-level
records into a temporal summary and reveal patterns that may not be apparent from
individual transactions.

\subsection{Overall Technical Interpretation}

From a technical perspective, the analysis demonstrates the suitability of distributed
computing for financial crime analytics. Operations such as:

\begin{itemize}
    \item Mapping transaction attributes
    \item Grouping records by payment format, bank, and account pairs
    \item Aggregating millions of records into summary statistics
\end{itemize}

were executed using PySpark's distributed architecture.

Instead of analysing more than 176 million transactions individually, distributed
reduction operations produced compact summaries that exposed meaningful laundering
patterns while maintaining computational efficiency. The results therefore demonstrate
both analytical value and the practical benefits of distributed data processing for large-scale financial datasets.